%%%%%%%%%%%%%%%%%
% This is an sample CV template created using altacv.cls
% (v1.3, 10 May 2020) written by LianTze Lim (liantze@gmail.com). Now compiles with pdfLaTeX, XeLaTeX and LuaLaTeX.
% This fork/modified version has been made by Nicolás Omar González Passerino (nicolas.passerino@gmail.com, 15 Oct 2020)
%
%% It may be distributed and/or modified under the
%% conditions of the LaTeX Project Public License, either version 1.3
%% of this license or (at your option) any later version.
%% The latest version of this license is in
%%    http://www.latex-project.org/lppl.txt
%% and version 1.3 or later is part of all distributions of LaTeX
%% version 2003/12/01 or later.
%%%%%%%%%%%%%%%%

%% If you need to pass whatever options to xcolor
\PassOptionsToPackage{dvipsnames}{xcolor}

%% If you are using \orcid or academicons
%% icons, make sure you have the academicons
%% option here, and compile with XeLaTeX
%% or LuaLaTeX.
% \documentclass[10pt,a4paper,academicons]{altacv}

%% Use the "normalphoto" option if you want a normal photo instead of cropped to a circle
% \documentclass[10pt,a4paper,normalphoto]{altacv}

\documentclass[10pt,a4paper,ragged2e,withhyper]{altacv}

% --- Chinese Support ---
\usepackage{xeCJK}
% "Microsoft YaHei" is standard on Windows and looks good for resumes (Sans-serif)
\setCJKmainfont[
  BoldFont={Microsoft YaHei}, 
  ItalicFont={Microsoft YaHei}
]{Microsoft YaHei}

%% GEMINI NOTE: should only use xelatex to compile

%% AltaCV uses the fontawesome5 and academicons fonts
%% and packages.
%% See http://texdoc.net/pkg/fontawesome5 and http://texdoc.net/pkg/academicons for full list of symbols. You MUST compile with XeLaTeX or LuaLaTeX if you want to use academicons.

% Change the page layout if you need to
\geometry{left=1.2cm,right=1.2cm,top=1cm,bottom=1cm,columnsep=0.75cm}

% The paracol package lets you typeset columns of text in parallel
\usepackage{paracol}
\usepackage{ragged2e}
\usepackage{hyperref}
\usepackage{tabto}

% Change the font if you want to, depending on whether
% you're using pdflatex or xelatex/lualatex
\ifxetexorluatex
  % If using xelatex or lualatex:
  \setmainfont{Roboto Slab}
  \setsansfont{Lato}
  \renewcommand{\familydefault}{\sfdefault}
\else
  % If using pdflatex:
  \usepackage[rm]{roboto}
  \usepackage[defaultsans]{lato}
  % \usepackage{sourcesanspro}
  \renewcommand{\familydefault}{\sfdefault}
\fi

% ----- LIGHT MODE -----
\definecolor{SlateGrey}{HTML}{2E2E2E}
\definecolor{LightGrey}{HTML}{666666}
\definecolor{PrimaryColor}{HTML}{001F5A}
\definecolor{SecondaryColor}{HTML}{0039AC}
\definecolor{ThirdColor}{HTML}{F3890B}
\definecolor{BackgroundColor}{HTML}{E2E2E2}
\colorlet{name}{PrimaryColor}
\colorlet{tagline}{PrimaryColor}
\colorlet{heading}{PrimaryColor}
\colorlet{headingrule}{ThirdColor}
\colorlet{subheading}{SecondaryColor}
\colorlet{accent}{SecondaryColor}
\colorlet{emphasis}{SlateGrey}
\colorlet{body}{LightGrey}
\pagecolor{BackgroundColor}   
% ----- DARK MODE -----
%\definecolor{BackgroundColor}{HTML}{242424}
%\definecolor{SlateGrey}{HTML}{6F6F6F}
%\definecolor{LightGrey}{HTML}{ABABAB}
%\definecolor{PrimaryColor}{HTML}{3F7FFF}
%\colorlet{name}{PrimaryColor}
%\colorlet{tagline}{PrimaryColor}
%\colorlet{heading}{PrimaryColor}
%\colorlet{headingrule}{PrimaryColor}
%\colorlet{subheading}{PrimaryColor}
%\colorlet{accent}{PrimaryColor}
%\colorlet{emphasis}{LightGrey}
%\colorlet{body}{LightGrey}
%\pagecolor{BackgroundColor}

% Change some fonts, if necessary
\renewcommand{\namefont}{\Huge\rmfamily\bfseries}
\renewcommand{\personalinfofont}{\small\bfseries}
\renewcommand{\cvsectionfont}{\large\rmfamily\bfseries}
\renewcommand{\cvsubsectionfont}{\large\bfseries}

% Change the bullets for itemize and rating marker
% for \cvskill if you want to
\renewcommand{\itemmarker}{{\small\textbullet}}
\renewcommand{\ratingmarker}{\faCircle}

%% sample.bib contains your publications
\addbibresource{sample.bib}
% \renewbibmacro{in:}{}

%% Gemini (reformat bib section)
\AtBeginBibliography{
  \small                          % 1. Set font to small (or try \footnotesize)
  \renewcommand{\itemmarker}{}    % 2. Remove the bullet point
  \setlength{\biblabelsep}{0.5em} % 3. Reduce the gap between [1] and the text
}

% --- Highlight Name in Bibliography ---
\renewcommand*{\mkbibnamegiven}[1]{%
  \ifboolexpr{ test {\ifstrequal{\namepartfamily}{Zhang}}
               and test {\ifstrequal{\namepartgiven}{Jiannan}} }
    {\textbf{#1}}{#1}%
}

\renewcommand*{\mkbibnamefamily}[1]{%
  \ifboolexpr{ test {\ifstrequal{\namepartfamily}{Zhang}}
               and test {\ifstrequal{\namepartgiven}{Jiannan}} }
    {\textbf{#1}}{#1}%
}

\begin{document}
    \name{张先生}
    \tagline{飞控系统工程师}
    %% You can add multiple photos on the left or right
    \photoL{4cm}{JZhang}
    
    \personalinfo{
        \email{jiannan.zhang@tum.de}\smallskip
        \phone{+49-176-630-81506}
        \location{德国 慕尼黑}\\
        \linkedin{jiannan-zhang}
        %\github{johnDoe}
        %\dev{johnDoe}
        %\homepage{nicolasomar.me}
        %\medium{nicolasomar}
        %% You MUST add the academicons option to \documentclass, then compile with LuaLaTeX or XeLaTeX, if you want to use \orcid or other academicons commands.
        % \orcid{0000-0000-0000-0000}
        %% You can add your own arbtrary detail with
        %% \printinfo{symbol}{detail}[optional hyperlink prefix]
        % \printinfo{\faPaw}{Hey ho!}[https://example.com/]
        %% Or you can declare your own field with
        %% \NewInfoFiled{fieldname}{symbol}[optional hyperlink prefix] and use it:
        % \NewInfoField{gitlab}{\faGitlab}[https://gitlab.com/]
        % \gitlab{your_id}
    }
    
    \makecvheader
    %% Depending on your tastes, you may want to make fonts of itemize environments slightly smaller
    % \AtBeginEnvironment{itemize}{\small}
    
    %% Set the left/right column width ratio to 6:4.
    \columnratio{0.27}

    % Start a 2-column paracol. Both the left and right columns will automatically
    % break across pages if things get too long.
    \begin{paracol}{2}
        % ----- STRENGTHS -----
        \cvsection{\large 核心优势}
            \cvtag{飞控系统}\\
            \cvtag{飞行物理与建模仿真}\\
            \cvtag{基于模型的系统工程}\\
            \cvtag{适航取证}\\
            % \medskip
            \cvtag{飞行测试}
            
        % ----- STRENGTHS -----
        
        % ----- LEARNING -----
        \cvsection{\large 软件技能}
            \cvtag{MATLAB/Simulink}\\
            \cvtag{\normalsize Control Design Toolbox}\\
            \cvtag{\normalsize Math Symbolic Toolbox}\\
            \cvtag{Git}
            \cvtag{C/C++}
        % ----- LEARNING -----
        
        % ----- LANGUAGES -----
        \cvsection{\large 语言能力}
            \cvlang{英语}{\small 商务}\\
            % \divider
            \medskip
            
            \cvlang{普通话}{\small 母语}\\
            \medskip
            % \divider
            
            
            \cvlang{德语}{\small 歌德B1}
            \smallskip
            %% Yeah I didn't spend too much time making all the
            %% spacing consistent... sorry. Use \smallskip, \medskip,
            %% \bigskip, \vpsace etc to make ajustments.
            
        % ----- LANGUAGES -----
            
        % ----- REFERENCES -----
        % \cvsection{References}
        %     \cvref{Ref 1}{ref-1}
        %     \divider
            
        %     \cvref{Ref 2}{ref-2}
        %     \divider
            
        %     \cvref{Ref 3}{ref-3}
        %     \smallskip
        % ----- REFERENCES -----
        
        % ----- MOST PROUD -----
        % \cvsection{Most Proud of}
        
        % \cvachievement{\faTrophy}{Fantastic Achievement}{and some details about it}\\
        % \divider
        % \cvachievement{\faHeartbeat}{Another achievement}{more details about it of course}\\
        % \divider
        % \cvachievement{\faHeartbeat}{Another achievement}{more details about it of course}
        % ----- MOST PROUD -----
        
        % \cvsection{A Day of My Life}
        
        % Adapted from @Jake's answer from http://tex.stackexchange.com/a/82729/226
        % \wheelchart{outer radius}{inner radius}{
        % comma-separated list of value/text width/color/detail}
        % \wheelchart{1.5cm}{0.5cm}{%
        %   6/8em/accent!30/{Sleep,\\beautiful sleep},
        %   3/8em/accent!40/Hopeful novelist by night,
        %   8/8em/accent!60/Daytime job,
        %   2/10em/accent/Sports and relaxation,
        %   5/6em/accent!20/Spending time with family
        % }
        
        % use ONLY \newpage if you want to force a page break for
        % ONLY the current column
        % \newpage
        
        %% Switch to the right column. This will now automatically move to the second
        %% page if the content is too long.
        \switchcolumn
        
        % ----- Objective -----
        \cvsection{个人简介}
        
            % Removing 'quote' restores normal text size and margins
            \begin{tabular}{@{}ll}
                \textbf{个人愿景:} & 实现eVTOL产品化. \\[0.2em]
                \textbf{工业经历:} & 10+ 年eVTOL飞控系统开发, 1张TC, 1000+小时调测经验.\\[0.2em]
                \textbf{学术经历:} & 10+篇核心期刊, 300+次被引用.\\ % [0.2em]
            \end{tabular}
            
            % \end{quote}
        % ----- Objective -----
        
        % ----- EXPERIENCE -----
        \cvsection{工作经历}
            \cvevent{科研助理}{| 慕尼黑工业大学}{06.2016 -- 至今}{德国慕尼黑}
            \begin{itemize}
                \item 负责控制分配算法开发。
                \item 负责控制构型优化设计。
                \item 飞行控制系统开发。
                \item 飞行动力学建模工具链开发。
            \end{itemize}
            % \smallskip
            \bigskip

            \cvevent{飞控系统负责人 }{| 峰飞航空科技}{08.2023 -- 07.2024}{中国上海}
            \begin{itemize}
                \item 负责制导、导航与控制 (GNC) 功能开发。
                \item 负责系统架构设计。
                \item 负责飞控系统适航取证 (TC) 申请工作。
            \end{itemize}
            % \smallskip
            \bigskip
            
            \cvevent{高级飞控工程师 }{| 峰飞航空科技}{07.2021 -- 08.2023}{中国上海}
            \begin{itemize}
                \item 负责制导与控制律开发。
                \item 负责飞行动力学建模开发。
                \item 参与系统架构设计。
                \item 参与适航认证任务（需求捕获与验证）。
                \item 参与飞行测试。
            \end{itemize}
            % \smallskip
            \bigskip
            
            \cvevent{飞控工程师 }{| 道通智能航空（欧洲）}{06.2016 -- 01.2020}{德国 Ismaning}
            \begin{itemize}
                \item 参与倾转旋翼无人机 (龙鱼) 飞控算法开发。
                \item 主导姿态环设计与控制分配。
                \item 参与飞行测试。
            \end{itemize}
            \bigskip

            \cvevent{软件开发工程师 }{| Phoenix-Wings GmbH}{06.2020 -- 07.2021}{德国 Ismaning}
            \begin{itemize}
                \item 负责配平与线性化工具开发。
                \item 负责仿真工具链开发。
                \item 参与飞行测试。
                \item 协助项目管理。
            \end{itemize}
            % \bigskip
            % \divider
        % ----- EXPERIENCE -----
    
    \end{paracol}
    \newpage
    % --- Set new, wider margins for Page 2 onwards ---
    % You can adjust '2.5cm' to whatever width you prefer (e.g., 2cm, 1in)
    \newgeometry{left=2.5cm, right=2.5cm, top=2cm, bottom=2cm}

        % ----- EDUCATION -----
        \cvsection{教育经历}
            \cvevent{博士在读 }{| 慕尼黑工业大学 (TUM)}{11.2017 -- 2025}{德国慕尼黑}
            \begin{itemize}
                \item \textbf{研究方向}: \small eVTOL 控制操纵性与分配优化
            \end{itemize}
            \bigskip
            % \divider
            
            \cvevent{航空航天工程（硕士） }{| 南洋理工大学 (NTU)}{08.2013 -- 02.2016}{新加坡}
            \begin{itemize}
                \item \textbf{CGPA}: 4.92 \small (CAP: 1.06) \normalsize  \hspace{1em} \textbf{主修课程}: 飞行控制，轻量化结构
            \end{itemize}
            \bigskip
            % \divider
            
            % \cvevent{Uncompleted}{| Beihang University}{09.2012 -- 07.2013}{Beijing, China}
            % \begin{itemize}
            %     \item First year courses complete.
            %     \item 50\% finished towards M.Sc. Aerospace Engineering.
            % \end{itemize}
            % \divider
            
            \cvevent{飞行器设计与工程（本科） }{| 北京理工大学 (BIT)}{08.2008 -- 07.2012}{中国北京}
            \begin{itemize}
                \item \textbf{GPA}: 85/100 \hspace{5em} \textbf{主修课程}: 飞行力学，空气动力学
            \end{itemize}
        % ----- EDUCATION -----
        
        % ----- PROJECTS -----
        % \cvsection{Projects}
        %     \cvevent{V280 }{\cvrepo{| \faGithub}{https://github.com/user/repo}\cvrepo{| \faGlobe}{https://repo-demo.com/}}{02.2020 -- Now}{}
        %     \begin{itemize}
        %         \item Item 1
        %         \item Item 2
        %     \end{itemize}
        %     \divider
        
        \renewcommand{\cvevent}[4]{%
          {\large\color{emphasis}#1}
          \ifstrequal{#2}{}{}{\large{\color{accent}#2}}
          \par\medskip\normalsize
          \ifstrequal{#3}{}{}{{\small\makebox[0.5\linewidth][l]{\color{accent}\faCalendar\color{emphasis}~#3}}}
          \ifstrequal{#4}{}{}{{\small\makebox[0.5\linewidth][l]{\color{accent}\faDev\color{emphasis}~#4}}}\par
          \medskip\normalsize%\faMapMarker
        }
        
        \cvsection{工业界项目}
            \cvevent{凯瑞欧 - 2000kg 复合翼eVTOL}{| 峰飞航空科技}{07.2021 -- 08.2024}{工业项目}
            \begin{itemize} \justifying
                \item \textbf{角色}: 建模、制导与控制律负责人
                % \item Major Requirements:
                % % \vspace{1mm}
                % \begin{itemize}
                % \setlength\itemsep{0.1ex}
                \item \textbf{贡献}: 控制律设计，软件实现，系统架构及适航认证任务 (FHA, PSSA 等)。
                \item \textbf{亮点}: 全球首个获得型号合格证 (TC) 的 2吨级 eVTOL。
                % \end{itemize}
            \end{itemize}
            \bigskip
        
            \cvevent{V280 - 280kg 复合翼eVTOL}{| TUM-FSD}{02.2020 -- 06.2021}{工业项目}
            \begin{itemize} \justifying
                \item \textbf{角色}: 控制分配开发负责人
                % \item Major Requirements:
                % % \vspace{1mm}
                % \begin{itemize}
                % \setlength\itemsep{0.1ex}
                \item \textbf{贡献}: 开发悬停、过渡及巡航阶段的控制分配算法，并通过 OEI (单发失效) 故障安全验证。
                \item \textbf{亮点}: 依据 DO-331 标准开发（使用标准库、模型顾问、测试框架等）。
            \end{itemize}
            \bigskip
            % \divider
            
            \cvevent{魔鬼鱼 - 30kg 复合翼无人机}{| Phoenix-Wings GmbH}{06.2020 -- 07.2021}{工业项目}
            \begin{itemize} \justifying
                \item \textbf{角色}: C++ 配平与线性化工具负责人
                \item \textbf{贡献}: 使用 C++ 实现配平与线性化工具，以及用于模型验证和后续设计任务的 MATLAB/Simulink 接口。
            \end{itemize}
            \bigskip

            \cvevent{龙鱼 - 7kg 倾转旋翼无人机}{| 道通智能航空（欧洲）}{06.2016 -- 01.2020}{工业项目}
            \begin{itemize} \justifying
                \item \textbf{角色}: 控制分配开发负责人
                \item \textbf{贡献}: 开发全包线非线性控制分配；支持 INDI 姿态与速度环设计；发表两篇一作论文。
                \item \textbf{亮点}: 作为成熟产品已产生数十亿销售额。
            \end{itemize}
            
            % \divider
        % ----- PROJECTS -----
        
        \cvsection{学术课题}
            \cvevent{广义控制分配优化与保护}{| TUM-FSD}{01.2025 -- 12.2025}{科研项目}
            \begin{itemize} \justifying
                \item \textbf{角色}: 算法开发负责人
                \item \textbf{贡献}: 提出一种广义p-范数优化目标的自适应方法；开发两种基于伪逆法的指令保护改进方法；发表一篇一作期刊和两篇二作期刊。
            \end{itemize}
            \bigskip

            \cvevent{eVTOL 控制构型优化}{| TUM-FSD}{05.2020 -- 12.2024}{科研项目}
            \begin{itemize} \justifying
                \item \textbf{角色}: 算法设计与软件实现负责人
                \item \textbf{贡献}: 开发了一种优化 eVTOL 动力升力旋翼拓扑的方法，最优构型在单发失效场景下具有极佳的鲁棒性；发表两篇一作期刊。
            \end{itemize}
            \bigskip
            % \divider 

            \cvevent{F-16 建模扩展与控制设计}{| TUM-FSD}{05.2016 -- 01.2020}{科研项目}
            \begin{itemize} \justifying
                \item \textbf{角色}: 建模与 INDI 基线控制器负责人
                \item \textbf{贡献}: 开发 INDI 闭环仿真框架，以及非线性最优控制分配和先进避障算法；发表两篇一作论文。
            \end{itemize}
            \bigskip
            % \divider    
        
        % ----- PROJECTS -----
        \cvsection{学术成果}
        \nocite{*} % The * means "include everything in the .bib file"
        \printbibliography[heading=none]

        % \begin{enumerate} \justifying \small 
        %     \item \emph{Attainable Moment Set Optimization to Support Configuration Design: A Required Moment Set Based Approach}
        %     \item \emph{Saturation Protection with Pseudo Control Hedging: A Control Allocation Perspective}
        %     \item \emph{Control Allocation Framework with SVD-based Protection for a Tilt-rotor VTOL Transition Air Vehicle}
        %     \item \emph{Modeling and Incremental Nonlinear Dynamic Inversion Control for a Highly Redundant Flight System}
        %     \item \emph{Control Allocation Framework for a Tilt-rotor Vertical Take-off and Landing Transition Aircraft}
        % \end{enumerate}
        
        
    
\end{document}
